\documentclass[Screen16to9,17pt]{foils}
\usepackage{zencurity-slides}
\selectlanguage{english}

% Software Security - Exam project FlyBook
% Contents

%
\begin{document}

\mytitlepage
{Exam Project}
{Software Security -- FlyBook}

\slide{FlyBook}

\begin{list2}
\item System that allows the user to search for flights and book them
\end{list2}

\hlkimage{10cm}{flybook.png}

\slide{The task}

\begin{list2}
\item You are asked to make a security assessment of the FlyBook system
\item Unfortunately, the documentation of the system is not available
\item The good news is that you are given access to the source code
\item You need to hand in a security assessment report
\end{list2}

\vskip 1cm
Source: \url{https://github.com/sebikea/flybook}

Note: if you have problems running the project work together with other groups!

Recommend cloning the project and making local changes - track changes



\slide{Report -- required contents}

Can be structured as you wish, but make sure it includes:
\begin{list1}
\item Description of the system and what kind of system it is, what is the context -- a half page that sets the scene, what kind of security does such a system warrant
\item Static analysis (code review)
  \begin{list2}
  \item Automated code review -- use at least two SAST tools\\
  Describe a little about the settings, maybe include some config in appendices
  \item Manual code review -- have you noticed something bad
  \end{list2}
\item Dynamic analysis (while code is running)
  \begin{list2}
  \item Automated analysis -- use at least one DAST tool -- could be dynamic audit/runtime or pentesting\\
  Describe a little about the settings, maybe include some config in appendices
  \item Manual pentesting (optional) -- try working with the application in Burp\\
  Document some of the findings -- use categories informational, low risk, medium risk or high risk
  \end{list2}
\end{list1}


\slide{Report -- required contents, cont.}

\begin{list1}
\item Security problems and recommendations (possible/suggested fixes)
  \begin{list2}
  \item Explain what should be changed and how (not detailed code rewrite or refactor)
  \item You are encouraged to provide references to CWE
  \end{list2}
\item Risk analysis to identify the most urgent vulnerabilities to be fixed\\
  (detailed risk analysis in the appendix)\\
  Priorities - easy fixes, critical fixes, where would you start!
\item Common report sections: introduction, conclusion,\\
  references/literature, and appendix (bilag) if needed
\end{list1}

\slide{Aspects to consider -- tools}

\begin{list2}
\item Which tools have you selected and used? Why?
\item What tools or testing approaches give you the best results and why?
\item If the tool allows it, tune it to have fewer false positives\\
  -- and explain in the report what you did
\item \textbf{Don't include all findings from the automated code review tools:}\\
  select the most relevant ones or summarize the results
\end{list2}

\slide{Aspects to consider -- findings}

\begin{list1}
\item What are the most serious vulnerabilities and why?
\item How is data transmitted between the different components?\\
  GET/POST requests, HTTPS?
\item Code review -- remember to look at:
  \begin{list2}
  \item Errors and exception handling
  \item Container security best practices
  \item Logging
\item You can also use Software Composition Analysis (SCA) tools (dependency checking)
\end{list2}
\end{list1}

\slide{Report -- formalities}

\begin{list2}
\item Has to be written in Danish (unless you have applied for permission)
\item Max. 10 A4 pages + 5 pages per student (approx.)
  \begin{list2}
  \item Group of 2 students: max. 10+10 pages
  \item Group of 3 students: max. 10+15 pages
  \item Group of 4 students: max. 10+20 pages
  \end{list2}
\item Do your best to work in groups of 2--4 students\\
  -- please try to avoid 1 person groups as it will be too much work!
\end{list2}

\slide{Exam project hand-in}

\begin{list2}
\item Through WiseFlow
\item Deadline: TBD, I will check with Dany
\item Remember to include the names of all group members in the report
\item Report should be in PDF
\item If you modify the code to fix the vulnerabilities (not mandatory)\\
you may include ZIP or reference your Git repository
\end{list2}

\slide{Exam}

\begin{list2}
\item Dates: TBD -- I think they are scheduled already!
\item Oral, individual exam
\item Language: Danish
\item 10 min presentation of the findings\\
  + 15 min questions based on the project and course curriculum
\item 7-point grading scale
\end{list2}

\slide{Presentation at the exam}

\begin{list1}
\item Not mandatory to use slides, but they can provide structure\\
  (try to prepare different slides per group member)
\item Max 10 min -- you can summarize the findings
\item You can include extra bits not covered in the report but seen in the course:
  \begin{list2}
  \item Software security touchpoints
  \item Security design principles
  \item \ldots
  \end{list2}
\item Make sure you understand at a high level the vulnerabilities\\
  you mention in your presentation -- what are the consequences of these
  \item You can do a demo of an attack -- not recommended, takes time!
\end{list1}

\end{document}
